% ---------------------------------------------------------------------
%  CHƯƠNG 5 — KẾT QUẢ VÀ ĐÁNH GIÁ
% ---------------------------------------------------------------------
\chapter{Kết quả và đánh giá}
\label{ch:ketqua}

\begin{muctieu}
Chương này đối chiếu kết quả với các mục tiêu đã đặt ra, sau đó trình bày
chi tiết sai số của từng phương pháp, kiểm định thống kê, độ phủ của nón bất
định, hai nhóm thí nghiệm bổ sung, kết quả kiểm thử phần mềm, phần bàn luận
và những hạn chế còn tồn tại.
\end{muctieu}

\section{Kết quả đạt được}

Bảng~\ref{tab:muctieu} đối chiếu từng mục tiêu đã nêu ở Chương~\ref{ch:modau}.
Bốn mục tiêu đầu đạt đầy đủ; mục tiêu triển khai đạt một phần vì mô hình tốt
nhất trong thí nghiệm chưa phải là mô hình đang phục vụ. Các mục sau trình bày
bằng chứng cho từng dòng của bảng.

\begin{table}[H]
\centering
\caption{Đối chiếu mục tiêu đồ án.}
\label{tab:muctieu}
\begin{tabularx}{\textwidth}{>{\bfseries\color{daunavy}}l X >{\leavevmode\color{daumaroon}\raggedright\arraybackslash}p{3.4cm}}
\toprule
\rowcolor{gray100}\bfseries\color{daunavy}Mục tiêu & \leavevmode\bfseries\color{daunavy}Kết quả & \leavevmode\bfseries\color{daumaroon}Mức độ \\
\midrule
MT1 & Quy trình dữ liệu 24.024 quan trắc, chia theo năm, scaler fit trên train & Đạt \\
MT2 & 7 phương pháp, 3 seed, bootstrap theo cơn bão & Đạt \\
MT3 & 72 giờ: 676~km (persistence) xuống 462~km (k-NN), giảm 32\% & Đạt \\
MT4 & Nón 80\% phủ 82--86\% vị trí thực trên test & Đạt \\
MT5 & API, bản đồ web, demo công khai kèm bão tương tự; mô hình đang phục vụ chưa dùng dòng dẫn đường & Đạt một phần \\
\bottomrule
\end{tabularx}
\end{table}

\section{Độ đo đánh giá}

Độ đo chính là sai số Haversine trung bình tại từng mốc dự báo, tính bằng
kilômét theo công thức~\eqref{eq:haversine}. Đây là khoảng cách trên mặt cầu
giữa vị trí gợi ý và vị trí thực, nên càng nhỏ càng tốt. Các mạng nơ-ron chạy
3 seed và được báo cáo dưới dạng trung bình $\pm$ độ lệch chuẩn; persistence,
CLIPER-proxy và k-NN là tất định nên chỉ có một giá trị.

Kết quả được tính trên hai tập. Bảng~\ref{tab:ketqua-all} dùng toàn bộ 84 cơn
bão kiểm thử (714 cửa sổ). Bảng~\ref{tab:ketqua-vn} chỉ dùng 27 cơn bão ảnh
hưởng Việt Nam (278 cửa sổ), là nhóm mà đề tài quan tâm nhất.

\begin{table}[H]
\centering
\caption{Sai số trung bình (km) trên toàn bộ tập kiểm thử 2021--2024.}
\label{tab:ketqua-all}
\begin{tabular*}{\textwidth}{@{\extracolsep{\fill}}>{\bfseries\color{daunavy}}l rrrrr}
\toprule
\rowcolor{gray100}\bfseries\color{daunavy}Phương pháp & \bfseries\color{daumaroon}6h & \bfseries\color{daumaroon}12h & \bfseries\color{daumaroon}24h & \bfseries\color{daumaroon}48h & \bfseries\color{daumaroon}72h \\
\midrule
Persistence       & 31 & 68 & 155 & 390 & 676 \\
CLIPER-proxy      & 30 & 65 & 151 & 388 & 677 \\
k-NN analog       & 32 & 64 & 133 & 290 & \textbf{462} \\
Gradient boosting & \textbf{29} & \textbf{60} & \textbf{129} & 292 & 476 \\
LSTM              & 33 $\pm$ 0 & 65 $\pm$ 1 & 133 $\pm$ 1 & \textbf{289} $\pm$ 1 & 464 $\pm$ 5 \\
Attention-LSTM    & 35 $\pm$ 3 & 68 $\pm$ 4 & 138 $\pm$ 5 & 297 $\pm$ 10 & 473 $\pm$ 9 \\
Transformer       & 33 $\pm$ 1 & 64 $\pm$ 1 & 133 $\pm$ 3 & 293 $\pm$ 3 & 468 $\pm$ 5 \\
\bottomrule
\end{tabular*}
\end{table}

\begin{table}[H]
\centering
\caption{Sai số trung bình (km) trên 27 cơn bão ảnh hưởng Việt Nam.}
\label{tab:ketqua-vn}
\begin{tabular*}{\textwidth}{@{\extracolsep{\fill}}>{\bfseries\color{daunavy}}l rrrrr}
\toprule
\rowcolor{gray100}\bfseries\color{daunavy}Phương pháp & \bfseries\color{daumaroon}6h & \bfseries\color{daumaroon}12h & \bfseries\color{daumaroon}24h & \bfseries\color{daumaroon}48h & \bfseries\color{daumaroon}72h \\
\midrule
Persistence       & 33 & 70 & 159 & 388 & 642 \\
CLIPER-proxy      & 32 & 68 & 159 & 414 & 710 \\
k-NN analog       & 33 & 65 & 134 & 290 & 476 \\
Gradient boosting & \textbf{30} & 61 & 131 & 295 & 491 \\
LSTM              & 32 $\pm$ 0 & 63 $\pm$ 0 & 133 $\pm$ 1 & 292 $\pm$ 6 & 468 $\pm$ 9 \\
Attention-LSTM    & 34 $\pm$ 3 & 66 $\pm$ 4 & 138 $\pm$ 4 & 298 $\pm$ 6 & 471 $\pm$ 3 \\
Transformer       & 32 $\pm$ 1 & \textbf{62} $\pm$ 1 & \textbf{130} $\pm$ 2 & \textbf{289} $\pm$ 3 & \textbf{463} $\pm$ 10 \\
\bottomrule
\end{tabular*}
\end{table}

Hai bảng cho cùng một bức tranh. Ở 6 và 12 giờ, mọi phương pháp chỉ cách
nhau vài kilômét. Từ 24 giờ, nhóm phương pháp học tách khỏi hai mốc so sánh,
và đến 72 giờ khoảng cách lên tới khoảng 200~km. Trong nhóm phương pháp học,
không có phương pháp nào dẫn đầu ở mọi mốc. Hình~\ref{fig:horizon} biểu diễn
xu hướng này: các đường gần như trùng nhau ở 6 giờ và tách dần từ 24 giờ trở
đi.

\begin{figure}[H]
\centering
\includegraphics[width=0.85\textwidth]{error_by_horizon.png}
\caption{Sai số trung bình theo thời hạn dự báo của các phương pháp.}
\label{fig:horizon}
\end{figure}

\subsection{Kiểm định thống kê}

Chênh lệch vài chục kilômét giữa hai phương pháp có thể chỉ do may rủi của
tập kiểm thử. Để kiểm tra, Bảng~\ref{tab:bootstrap} so sánh tại 72 giờ bằng
bootstrap theo cơn bão: 50 cơn bão kiểm thử có dự báo 72 giờ được lấy mẫu lại
5.000 lần. Đơn vị lấy mẫu là cả cơn bão chứ không phải từng cửa sổ, vì các
cửa sổ của cùng một cơn bão không độc lập với nhau. Giá trị âm nghĩa là tốt
hơn mốc tham chiếu.

\begin{table}[H]
\centering
\caption{Chênh lệch sai số 72 giờ (km) và khoảng tin cậy 95\%.}
\label{tab:bootstrap}
\begin{tabular*}{\textwidth}{@{\extracolsep{\fill}}>{\bfseries\color{daunavy}}l cc}
\toprule
\rowcolor{gray100}\bfseries\color{daunavy}Phương pháp & \bfseries\color{daumaroon}So với persistence & \bfseries\color{daumaroon}So với LSTM \\
\midrule
CLIPER-proxy      & $-15$ \ [$-101$; $+65$]  & $+188$ \ [$+113$; $+256$] \\
k-NN analog       & $-201$ \ [$-268$; $-131$] & $+3$ \ [$-29$; $+34$] \\
Gradient boosting & $-191$ \ [$-252$; $-126$] & $+12$ \ [$-20$; $+48$] \\
LSTM              & $-203$ \ [$-262$; $-145$] & --- \\
Attention-LSTM    & $-173$ \ [$-229$; $-116$] & $+30$ \ [$+12$; $+50$] \\
Transformer       & $-200$ \ [$-257$; $-142$] & $+3$ \ [$-16$; $+22$] \\
\bottomrule
\end{tabular*}
\end{table}

Bảng cho ba kết luận. Mọi phương pháp học đều tốt hơn persistence 170--200~km
với khoảng tin cậy không chứa 0, nên mục tiêu MT3 đạt một cách chắc chắn.
Giữa LSTM, Transformer, k-NN và gradient boosting không có khác biệt có ý
nghĩa thống kê. Riêng Attention-LSTM kém hơn LSTM có ý nghĩa, và CLIPER-proxy
không khác persistence.

\subsection{Nón bất định}

Bán kính nón, lấy bằng phân vị 80\% của sai số trên validation, là 52, 97,
210, 455 và 714~km tại 6, 12, 24, 48 và 72 giờ. Trên tập kiểm thử, tỉ lệ vị
trí thực nằm trong nón là 82--86\% ở các mốc. Con số này sát mức danh nghĩa
80\% và hơi cao hơn, tức nón được hiệu chỉnh tốt và nghiêng về phía thận
trọng. Cần lưu ý nón ở 72 giờ rộng hơn 700~km: nó trung thực về mức độ bất
định, nhưng cũng cho thấy gợi ý ở hạn dài chỉ nên dùng để định hướng.

\subsection{Thí nghiệm loại bỏ: dự báo độ dịch chuyển}
\label{sec:ablation}

Chương~\ref{ch:thietke} đã lập luận rằng nên cho mạng học độ dịch chuyển thay
vì tọa độ tuyệt đối. Thí nghiệm này kiểm chứng lập luận đó: giữ nguyên kiến
trúc LSTM, dữ liệu và seed 0, chỉ đổi mục tiêu học (Bảng~\ref{tab:ablation}).

\begin{table}[H]
\centering
\caption{Ảnh hưởng của cách biểu diễn đầu ra lên sai số LSTM (km).}
\label{tab:ablation}
\begin{tabular*}{\textwidth}{@{\extracolsep{\fill}}>{\bfseries\color{daunavy}}l rrrrr}
\toprule
\rowcolor{gray100}\bfseries\color{daunavy}Mục tiêu học & \bfseries\color{daumaroon}6h & \bfseries\color{daumaroon}12h & \bfseries\color{daumaroon}24h & \bfseries\color{daumaroon}48h & \bfseries\color{daumaroon}72h \\
\midrule
Tọa độ tuyệt đối         & 115 & 121 & 164 & 305 & 473 \\
Độ dịch chuyển (đang dùng) & \textbf{33} & \textbf{65} & \textbf{133} & \textbf{289} & \textbf{464} \\
\bottomrule
\end{tabular*}
\end{table}

Khác biệt rõ nhất ở hạn ngắn. Học tọa độ tuyệt đối làm sai số 6 giờ tăng
khoảng 3,5 lần, kém cả persistence (31~km), và cần nhiều epoch hơn mới hội tụ
(dừng ở epoch 42 thay vì 13--17). Ở 72 giờ hai cách gần bằng nhau, vì khi đó
sai số do hướng đi tương lai quyết định chứ không do vị trí xuất phát. Kết
quả này xác nhận lựa chọn thiết kế: phần mạng cần học là chuyển động, không
phải vị trí hiện tại.

\subsection{Tối ưu huấn luyện}
\label{sec:steering}

Sau khi các phương pháp học cho kết quả ngang nhau, câu hỏi tiếp theo là còn
có thể giảm sai số bằng cách nào. Hai hướng được thử: thay đổi cách huấn
luyện trên cùng thông tin đầu vào, và bổ sung thông tin đầu vào mới. Các thí
nghiệm này đánh giá trên cùng các cửa sổ của bão vùng Việt Nam có đủ 48 giờ
lịch sử (435 cửa sổ validation, 526 cửa sổ test). Mọi lựa chọn dựa trên
validation; test chạy một lần.

\textbf{(a) Chỉ dùng thông tin quỹ đạo: đã chạm trần.} LSTM gốc có sai số
trung bình trên validation là 201~km. Không thay đổi nào cải thiện được con
số này: dùng dữ liệu toàn lưu vực (212~km), lịch sử 48 giờ (212~km), hàm mất
mát theo km (203~km), mạng lớn hoặc nhỏ hơn (206--207~km), Transformer
(207~km). Mô hình kết hợp LSTM và Transformer trên test chỉ giảm 8~km tại 72
giờ, khoảng tin cậy [$-16$; $+2$], không có ý nghĩa thống kê.

\textbf{(b) Thêm dòng dẫn đường khí quyển: cải thiện có ý nghĩa.} Hướng thứ
hai bổ sung 11 đặc trưng từ NCEP/NCAR R1 tại đúng thời điểm quan trắc
(Bảng~\ref{tab:dactrung}). Kết quả ở Bảng~\ref{tab:steering}.

\begin{table}[H]
\centering
\caption{Hiệu quả của đặc trưng dòng dẫn đường trên tập kiểm thử (km, 3 seed).}
\label{tab:steering}
\begin{tabular*}{\textwidth}{@{\extracolsep{\fill}}>{\bfseries\color{daunavy}}l rrrrrr}
\toprule
\rowcolor{gray100}\bfseries\color{daunavy}Cấu hình & \bfseries\color{daumaroon}6h & \bfseries\color{daumaroon}12h & \bfseries\color{daumaroon}24h & \bfseries\color{daumaroon}48h & \bfseries\color{daumaroon}72h & \bfseries\color{daumaroon}TB \\
\midrule
LSTM chỉ quỹ đạo (gốc)        & 32 & 64 & 131 & 286 & 469 & 197 \\
LSTM + dòng dẫn đường         & 30 & 58 & 118 & 270 & 465 & 188 \\
Transformer + dòng dẫn đường  & 30 & 57 & 115 & 265 & 450 & 183 \\
Kết hợp LSTM + Transformer    & \textbf{28} & \textbf{54} & \textbf{112} & \textbf{259} & \textbf{443} & \textbf{179} \\
\bottomrule
\end{tabular*}
\end{table}

Mô hình kết hợp (chọn trên validation) giảm 35~km tại 72 giờ so với mô hình
gốc, khoảng tin cậy 95\% [$-62$; $-9$] trên 44 cơn bão, có ý nghĩa thống kê;
tại 24 giờ giảm khoảng 15\%, trung bình các mốc giảm khoảng 9\%.

\section{Kiểm thử phần mềm}

Ngoài đánh giá mô hình, phần mềm được kiểm tra bằng bộ kiểm thử tự động viết
với \code{pytest}. Bộ kiểm thử bao phủ các khâu nạp dữ liệu, hợp đồng đặc
trưng, vòng huấn luyện, đóng gói và nạp gói mô hình, suy luận, API, lưu trữ,
bộ nhớ đệm và tiến trình cập nhật. Ở phiên bản bàn giao, 269 ca kiểm thử chạy
đạt và 1 ca được bỏ qua vì cần một máy chủ PostGIS bên ngoài.

Dữ liệu dùng trong kiểm thử phần mềm là dữ liệu tổng hợp. Các ca này chỉ
chứng minh phần mềm hoạt động đúng hợp đồng, không phải bằng chứng về độ
chính xác dự báo; độ chính xác chỉ được kết luận từ các bảng ở mục trên.

\section{Bàn luận}

\textbf{Học từ dữ liệu giúp rõ ở hạn dài.} Ở 6 giờ persistence gần như tối ưu,
vì trong thời gian ngắn như vậy bão chưa kịp đổi hướng. Từ 24 giờ trở đi, các
phương pháp học giảm sai số khoảng 14\% (24 giờ) đến 30\% (72 giờ). Lý do là
chúng học được các kiểu quỹ đạo điển hình, chẳng hạn đi về phía tây rồi
chuyển hướng lên phía bắc, điều mà phép ngoại suy thẳng không thể biết.

\textbf{Mô hình sâu không vượt trội.} Với khoảng 7.800 cửa sổ huấn luyện và 4
bước đầu vào, k-NN analog và gradient boosting đạt cùng mức với LSTM và
Transformer. Mạng nơ-ron overfit nhanh, epoch tốt nhất chỉ là 2--9. Thí
nghiệm (a) cho thấy đổi kiến trúc hay siêu tham số không giúp gì, còn thí
nghiệm (b) cho thấy thêm thông tin khí quyển thì có. Từ đó rút ra kết luận
chính của đồ án: nút thắt là thông tin đầu vào, không phải sức chứa của mô
hình. Điều này hợp với vật lý, vì bão bị dòng dẫn đường đẩy đi, và dữ liệu
chỉ có quỹ đạo không chứa thông tin này.

\textbf{Sai số theo cường độ.} Tại 72 giờ, áp thấp nhiệt đới là nhóm khó nhất
với persistence (800~km), còn bão rất mạnh cấp 12--13 là nhóm dễ nhất (khoảng
400~km với các phương pháp học). Một giả thuyết cần kiểm chứng thêm là hướng
đi của áp thấp kém ổn định hơn.

\textbf{Nghiên cứu trường hợp.} Sai số trung bình che mất những ca khó. Vì
vậy Bảng~\ref{tab:case} chọn thời điểm phát dự báo sao cho vị trí sau 72 giờ
gần bờ biển Việt Nam nhất, tức trường hợp bão sắp đổ bộ và gợi ý có giá trị
sử dụng cao nhất.

\begin{table}[H]
\centering
\caption{Sai số 72 giờ (km) khi bão sắp đổ bộ Việt Nam.}
\label{tab:case}
\small
\begin{tabular*}{\textwidth}{@{\extracolsep{\fill}}>{\bfseries\color{daunavy}}l rrrrr}
\toprule
\rowcolor{gray100}\bfseries\color{daunavy}Bão (thời điểm phát) & \bfseries\color{daumaroon}k-NN & \bfseries\color{daumaroon}GB & \bfseries\color{daumaroon}LSTM & \bfseries\color{daumaroon}Transf. & \bfseries\color{daumaroon}Persist. \\
\midrule
YAGI 2024 (06Z 04/09)    & 419 & \textbf{407} & 541 & 678 & 701 \\
SAOLA 2023 (12Z 31/08)   & \textbf{365} & 538 & 534 & 514 & 534 \\
YINXING 2024 (06Z 09/11) & 575 & \textbf{269} & 700 & 550 & 958 \\
CONSON 2021 (00Z 10/09)  & 558 & 581 & 625 & 702 & 987 \\
\bottomrule
\end{tabular*}
\end{table}

Với YAGI (Hình~\ref{fig:yagi}), phần lớn mô hình gợi ý bão chuyển hướng lên
phía bắc, kiểu quỹ đạo phổ biến trong dữ liệu, trong khi YAGI đi thẳng về phía
tây vào miền Bắc Việt Nam. Gradient boosting (407~km) và k-NN analog (419~km)
gần đúng nhất; k-NN dựa trên các bão tương tự như WUKONG 2000, DAMREY 2005. Với SAOLA, k-NN tốt nhất (365~km so với
514--579~km của các mạng nơ-ron). Với CONSON, CLIPER-proxy lại tốt nhất
(101~km); đây là trường hợp đơn lẻ, không đủ để kết luận. Các ví dụ này cho thấy
giá trị của k-NN không chỉ ở độ chính xác mà ở chỗ người dùng thấy được gợi ý
dựa trên những cơn bão nào.

\section{Hạn chế}

\begin{luuy}
\begin{itemize}[leftmargin=1.4em, itemsep=2pt]
  \item Mô hình đang phục vụ trên website và API vẫn là LSTM chỉ dùng quỹ đạo;
        mô hình có dòng dẫn đường cần dữ liệu khí quyển tại thời điểm dự báo
        (ví dụ phân tích GFS) mới đưa vào vận hành được.
  \item YAGI 2024 vẫn sai 620--670~km tại 72 giờ với mọi mô hình có dòng dẫn
        đường; lưới 2,5° của NCEP/NCAR và các đặc trưng tóm tắt chưa đủ để nhận
        ra áp cao cận nhiệt mạnh trong ca này.
  \item Nón bất định là đường tròn đẳng hướng, trong khi sai số dọc và ngang
        quỹ đạo thực tế khác nhau.
  \item Kết quả chỉ trên một cách chia theo năm; chưa kiểm định chéo theo mùa.
  \item Dữ liệu bão đang hoạt động từ IBTrACS ACTIVE trễ vài ngày.
  \item Hệ thống chưa được nghiệm thu trong môi trường vận hành thật: chưa có
        bằng chứng về tải, sao lưu và khôi phục trên máy chủ đích.
\end{itemize}
\end{luuy}
